\slajdid{09_1-s048}
\begin{frame}[label=09_1-s048]
\gusttekst\setlength{\parskip}{2pt}
\begin{columns}[c,onlytextwidth]\column{.35\textwidth}Karakteristika prenosa\column{.63\textwidth}\centering\begin{tikzpicture}[x=.67cm,y=.67cm,font=\sitneoznake,>=Latex,baseline=(current bounding box.center)]
\draw[->](-.13,0)--(3.45,0)node[right]{$V_I$};\draw[->](0,-.13)--(0,2.9)node[above]{$V_O$};\node[left]at(0,2.4){$V_{OH}$};\node[left]at(0,.13){$V_{OL}$};
\draw(0,2.4)--(.28,2.4)..controls(.6,2.4)and(.65,2.05)..(.83,1.65)..controls(1.15,.7)and(1.3,.4)..(1.6,.3)--(2.85,.13);
\draw[dashed](.58,-.1)node[below]{$V_{IL}$}--(.58,2.26);\draw[dashed](1.54,-.1)node[below]{$V_{IH}$}--(1.54,.36);\node[below]at(2.85,0){$V_{DD}$};\draw[dashed](0,.13)--(2.85,.13);

\draw(.4,2.44)--(.78,2.06);\node[right]at(.73,2.43){$a{=}{-}1$};\draw(1.31,.59)--(1.77,.13);\node[right]at(1.66,.9){$a{=}{-}1$};
\end{tikzpicture}

\end{columns}\medskip
Bez želje da tražimo sada „tačne“ vrednosti ono što je evidentno jeste da izrazi liče na one koje smo dobilo kada smo analizirali invertor sa pasivnim opterećenjem, kao i da sve ove parametre možemo da podešavamo odnosom $\frac{k_{n,D}}{k_{n,L}}$ odnosno direktno odnosima $\frac{W_{n,D}}{W_{n,L}}$.\par\medskip
Rezultati koji su dobijeni ne bi trebalo da su iznenađujući pošto se tranzistor $T_L$ ponaša kao dinamička otpornost. Kada je ulazni napon mali, odnosno izlazni napon visok, i napon $V_{GS,L}$ je mali, dinamička otpornost tranzistora $T_L$ $r_{DSn,L}\approx\frac1{k_{n,L}(V_{GS,L}-V_{Tn,L})}$ je velika i obezbeđuje veliko pojačanje u prelaznoj zoni, dobru karakteristiku. Kada je ulazni napon visok, odnosno izlazni napon mali, i napon $V_{GS,L}$ je veliki, dinamička otpornost tranzistora $T_L$ je mala i obezbeđuje brzo punjenje kapacitivnosti koje opterećuju izlaz prilikom prelaska sa logičke nule na logičku jedinicu. Normalno ne može baš sve da se dobije idealno, pošto će ova mala dinamička otpornost uticati na napon logičke nule (povećavaće), strujni kapacitet logičke nule (smanjivaće) itd.
\end{frame}
